%==============================================================================
% Sjabloon projectvoorstel graduaatsproject
%==============================================================================
\documentclass{hogent-article}
\addbibresource{voorstel.bib}
%%=============================================================================
%% Projectgegevens 
%%=============================================================================

%% ---------- Jij ------------------------------------------------------------
\newcommand{\studentnaam}{Sarah Turner Burkitt}
\newcommand{\studentemail}{sarah.turnerburkitt@student.hogent.be}

%% Wordt gebruikt als bestandsnaam van je PDF's, bijvoorbeeld
%% Familienaam-Voornaam-rapport.pdf. Geen spaties of leestekens gebruiken.
\newcommand{\bestandsnaam}{TurnerBurkitt-Sarah}

%% ---------- Je project -----------------------------------------------------
\newcommand{\projecttitel}{Ontwikkeling van een visuele e-maileditor}
\newcommand{\projectondertitel}{Proof of Concept voor aanpasbare e-mailtemplates binnen Vesalius.ai}
\newcommand{\projectpitch}{%
    Dit graduaatsproject ontwikkelt een proof of concept van een
    WYSIWYG-e-maileditor waarmee gebruikers automatisch gegenereerde
    e-mailtemplates kunnen aanpassen, opslaan, vooraf bekijken en testen.
    De oplossing wordt gevalideerd op veilige HTML-verwerking en compatibiliteit
    met gangbare e-mailclients als basis voor een mogelijke latere integratie
    in Vesalius.ai.%
}

%% ---------- Waar je het doet ------------------------------------------------
\newcommand{\bedrijfsnaam}{Vesalius.ai}

%% Je mentor op de werkvloer (het bedrijf)
\newcommand{\mentornaam}{Bart D'Hoore}
\newcommand{\mentoremail}{Bart.D'Hoore@endare.be}

%% Je begeleider van de opleiding (HOGENT)
\newcommand{\begeleidernaam}{Luc Vervoort}
\newcommand{\begeleideremail}{luc.vervoort@hogent.be}

%% ---------- Je repository ---------------------------------------------------
\newcommand{\repolink}{https://github.com/STB95/Graduaatsproject_TurnerBurkittSarah_26-27}
%% op je poster: hij is een beoordeeld onderdeel van je oplevering.
\newcommand{\repolinkdocent}{https://les.lucvervoort.org/sarahturnerburkitt/graduaatsproject}

%% ---------- Vast voor dit opleidingsonderdeel -------------------------------
%% Deze twee hoef je normaal niet aan te passen.
\newcommand{\opleidingnaam}{Graduaat in het Programmeren}
\newcommand{\academiejaarwaarde}{2026-2027}

%% Zoekwoorden voor het voorstel en de poster (3 tot 5)
\newcommand{\projectkeywords}{Vesalius.ai, Angular, WYSIWYG e-maileditor}


%% Opleiding, vak en soort opdracht
\studyprogramme{Graduaat Programmeren}
\course{Graduaatsproject}
\assignmenttype{Projectvoorstel}

\academicyear{2026-2027 1e semester}

\title{Ontwerp en proof of concept van een WYSIWYG-e-maileditor voor het Vesalius-platform}

\author{Sarah Turner Burkitt}
\email{turnerb_sarah@hotmail.com}

%% Je mentor op de werkvloer
\supervisor[Mentor]{Bart D'Hoore (Endare, Bart.D'Hoore@endare.com)}

%% De link naar je repository
\projectrepo{https://github.com/STB95/Graduaatsproject_TurnerBurkittSarah_26-27}

\keywords{Angular, e-mailcompatibility, HTML-rendering, e-mailtemplates, WYSIWYG e-maileditor, Vesalius.ai}

\begin{document}

\begin{abstract}

\noindent\textbf{Probleemstelling}

Binnen Vesalius.ai zijn de mogelijkheden om automatisch verstuurde
e-mails, zoals afspraakbevestigingen, herinneringen en annulaties,
aan te passen beperkt. Hierdoor kunnen artsen en organisaties de
inhoud en opmaak onvoldoende afstemmen op hun communicatiebehoeften.

\medskip
\noindent\textbf{Voorgestelde oplossing}

Dit graduaatsproject omvat de ontwikkeling van een zelfstandige
proof of concept (PoC) van een WYSIWYG-editor waarmee gebruikers
e-mailtemplates visueel kunnen bewerken, opslaan, vooraf bekijken
en als testmail versturen. De oplossing ondersteunt dynamische
placeholders en meertalige templates.

\medskip
\noindent\textbf{Technische aanpak}

De editor wordt ontwikkeld met Angular en TypeScript, aangevuld
met een mock-API op basis van Node.js en Express. Voor de tijdelijke
testopstelling wordt een geschikte opslagoplossing gekozen.
De focus ligt op veilige HTML-verwerking en een correcte weergave
in Gmail, Hotmail/Outlook.com en Apple Mail op macOS en iPhone.
Docker, Mailpit en geautomatiseerde tests ondersteunen de validatie.
De ontwikkeling gebeurt bij voorkeur binnen de bestaande
Bitbucket-omgeving van Vesalius.ai.

\medskip
\noindent\textbf{Verwachte oplevering}

De beoogde resultaten zijn:
\begin{itemize}
    \item Een werkende e-maileditor met voorbeeldweergave en
    testmailfunctionaliteit.
    \item Een mock-API voor het laden, opslaan en verwerken van templates.
    \item Een reproduceerbare testomgeving met fictieve gegevens
    en geautomatiseerde tests.
    \item Testresultaten van controles in verschillende e-mailclients.
    \item Technische documentatie en, indien de beschikbare tijd
    dit toelaat, een voorstel voor toekomstige integratie in Vesalius.ai.
\end{itemize}

Het project biedt Vesalius.ai inzicht in de technische haalbaarheid,
beperkingen en mogelijke meerwaarde van een visuele e-maileditor
als basis voor een eventuele toekomstige integratie.

\end{abstract}

%%=============================================================================
%% De inhoud van je projectvoorstel
%%=============================================================================

%---------- 1. Context en probleemstelling ------------------------------------
\section{Context en probleemstelling}%
\label{vst:context}

\textbf{Context}
\begin{itemize}
    \item Een digitaal platform dat artsen ondersteunt 
bij consultaties en administratieve taken. 
\end{itemize}

\textbf{Probleemstelling}
\begin{itemize}
    \item De inhoud en opmaak van verstuurde e-mails
zijn onvoldoende aanpasbaar. 
\end{itemize}
Gevolg: communicatie tussen arts en patiënt 
voldoen niet aan beide partijen hun specifieke behoeftes. 

%---------- 2. Doelstelling ----------------------------------------------------
\section{Doelstelling}%
\label{vst:doelstelling}
\begin{itemize}
    \item Een proof of concept (PoC) ontwikkelen van een WYSIWYG-e-maileditor
    \item Gebruikers kunnen templates* aanpassen, opslaan, vooraf bekijken en als testmail versturen 
\end{itemize}
templates*: afspraakbevestiging, afspraakherinnering en afspraakannulering
 
\newpage
%---------- 3. Technologieën --------------------------------------------------
\section{Technologieën}%
\label{vst:technologieen}
\textbf{Frontend}
\begin{itemize}
    \item Angular 22 (opbouw van de interactieve UI)
    \item TypeScript(voor typed en onderhoud code)
    \item Quill 2  (basis voor de visuele WYSIWYG-editor)
\end{itemize}

\medskip
\textbf{Mock-API en opslag}
\begin{itemize}
    \item Node.js (uitvoering van de servercode)
    \item TypeScript(getypeerde serverlogica)
    \item Express(opzetten van API-endpoints voor het laden en opslaan van templates)
\end{itemize}

\medskip
\textbf{HTML-verwerking en e-mailrendering}
\begin{itemize}
    \item Handlebars(templateverwerking en het behoud van de vaste layoutschil)
    \item \texttt{sanitize-html}(om onveilige HTML te verwijderen)
    \item Placeholdervervanging met escaping (om dynamische waarden gecontroleerd en veilig in te voegen)
    \item Juice (om CSS inline te plaatsen voor een betere ondersteuning door e-mailclients)
    \item \texttt{html-to-text} (om naast HTML ook een tekstuele e-mailversie te genereren)
\end{itemize}

\medskip
\textbf{Testomgeving}
\begin{itemize}
    \item Docker (reproduceerbare lokale testomgeving)
    \item Mailpit (om uitgaande testmails lokaal op te vangen en te controleren)
\end{itemize}

\medskip
\textbf{Geautomatiseerde tests}
\begin{itemize}
    \item Vitest (unit- en integratietests)
    \item Playwright (browsercontroles en screenshotvergelijkingen)
\end{itemize}

\medskip
\textbf{Versiebeheer}
\begin{itemize}
    \item Bitbucket (versiebeheer, samenwerking en het bijhouden van codewijzigingen)
\end{itemize}


%---------- 4. Afbakening ------------------------------------------------------
\section{Afbakening}%
\label{vst:afbakening}
\textbf{Out-Of-Scope (OOS)}
\begin{itemize}
    \item Integratie in productieplatform
    \item Plannen of versturen van productie e-mails
    \item Keycloak e-mails met echte patiëntgegevens en AI-functionaliteiten
\end{itemize}
De bestaande layoutschil met logo, merkkleuren en footer blijft behouden. 
Een mogelijk integratievoorstel kan als de kernfunctionaliteiten en belangrijkste tests correct en tijdig afgerond zijn.

%---------- 5. Verwachte meerwaarde -------------------------------------------
\section{Verwachte meerwaarde}%
\label{vst:meerwaarde}

\textbf{Inzichten in}
\begin{itemize}
    \item Technische haalbaarheid, beperkingen en veiligheidsaspecten van visuele e-mailbewerking. 
\end{itemize}

\medskip  
\textbf{Deliverables}
\begin{itemize}
    \item Een werkende demonstratie
    \item Een reproduceerbare testopstelling 
    \item Geautomatiseerde tests en testresultaten van verschillende e-mailclients
\end{itemize}


%---------- 6. Planning en werkbelasting --------------------------------------
\section{Planning en werkbelasting}%
\label{vst:planning}

\begin{center}
\small
\begin{tabular}{p{0.36\linewidth}p{0.47\linewidth}r}
\toprule
\textbf{Fase} & \textbf{Resultaat} & \textbf{Uren} \\
\midrule
Analyse & Scope, vereisten en testcriteria & 12 \\
Technologiekeuze en opzet & Technologiekeuzes en ontwikkelomgeving & 10 \\
Ontwerp & Architectuur, editorconcept en testmatrix & 13 \\
Realisatie & Editor, templates, preview en testmails & 40 \\
Testen en validatie & Geautomatiseerde tests en clientcontroles & 15 \\
Rapport, poster en demo & Eindrapport en demonstratie & 10 \\
\midrule
\textbf{Totaal} & & \textbf{100} \\
\bottomrule
\end{tabular}
\end{center}

%---------- 7. Risico's --------------------------------------------------------
\section{Risico's}%
\label{vst:risicos}

\textbf{Beperkte werkbelasting}
\begin{itemize}
    \item Niet alle uitbreidingen worden tijdig gerealiseerd
\end{itemize}

\textbf{Verschillen tussen e-mailclients}
\begin{itemize}
    \item Gegenereerde inhoud en opmaak wordt verschillend weergegeven (Outlook)
    \item CSS inline 
    \item Output tests in Gmail, Outlook en Apple
    \item Outlook Desktop wordt meegenomen voor zover de beschikbare testomgeving het toelaat
\end{itemize}

\textbf{Verschillen met de bestaande backend}
\begin{itemize}
    \item De mock-API kan afwijken van de bestaande Laravel-backend van Vesalius.ai
    \item Verschillen in architectuur en gegevensverwerking
\end{itemize}



\printbibliography[heading=bibintoc]

\end{document}
